\documentclass[12pt]{article}
\usepackage[utf8]{inputenc}
\usepackage{mathptmx} % Times New Roman
\usepackage{amsmath}
\usepackage{setspace}
\doublespacing
\usepackage[margin=1in]{geometry}
\usepackage{titlesec}
\usepackage{enumitem}

% AMA citation style (superscript numbers)
\usepackage[numbers,super]{natbib}
\bibliographystyle{ama}

\title{\textbf{Anchoring Large Language Models to Empirical Behavioral Evidence}\\[0.5em]
\large A Hybrid Approach for DCE-Grounded Narrative Policy Simulation}
\author{}
\date{}

\begin{document}

\maketitle

\begin{abstract}
\noindent\textbf{Objective:} To develop empirical-frontier regularization (EFR), an inference-time procedure anchoring LLM outputs to a discrete choice experiment (DCE) reference surface.

\noindent\textbf{Methods:} A Wuhan vaccination DCE ($N=1{,}027$) was combined with LLM simulation on a 64-state policy grid. EFR computes $\pi_{\text{EFR}} = \lambda\pi_{\text{LLM}} + (1-\lambda)P_{\text{emp}}$ with $\lambda=0.25$.

\noindent\textbf{Results:} Independent local replication validation on six administered DCE tasks ($N=205$ respondents) showed multinomial log loss: Pure-DCE 0.911, EFR 0.923, Raw LLM 1.020, Uniform 1.099. EFR improved over Raw LLM (9.5\% log loss reduction) and approached Pure-DCE (within 1.3\%). On a 64-state grid, waiting-time monotonicity violations fell from 10.4--62.5\% to 7.3--30.2\% across three LLMs under EFR.

\noindent\textbf{Limitations:} All respondents completed identical tasks; validation evaluates generalization to unseen respondents, not unseen tasks. Out-of-design scenarios lacked human ground truth.

\noindent\textbf{Conclusion:} EFR is a constrained narrative-extension mechanism for an existing DCE, not a replacement for the underlying model or validation of LLMs as standalone behavioral simulators.
\end{abstract}

\vspace{1em}
\noindent\textbf{Keywords:} empirical-frontier regularization; large language models; discrete choice experiment; health technology assessment

\vspace{1em}
\noindent\textbf{Article Type:} Methodological Article

\vspace{2em}
\section*{Highlights (120 words maximum)}
\begin{enumerate}[leftmargin=*]
    \item EFR is an inference-time regularization layer that constrains LLM outputs to a DCE-estimated reference without retraining.
    \item On held-out respondent validation, EFR recovers most of the DCE model's predictive performance while retaining a bounded LLM component.
    \item EFR improves structural consistency (monotonicity, rank order) but does not establish behavioral validity for out-of-design narrative scenarios.
\end{enumerate}

\newpage
\section{Introduction}
\label{sec:intro}

Large language models (LLMs) are increasingly used as simulated decision-makers for healthcare policy. Their appeal is interpretability of rich natural-language scenarios that structured models cannot directly process. However, linguistic fluency does not guarantee behavioral reliability: unconstrained LLMs may violate basic preference relationships implied by observed human behavior.

We propose empirical-frontier regularization (EFR), an inference-time procedure anchoring LLM outputs to a DCE-estimated reference surface without retraining. The method blends raw LLM probabilities with an empirical choice frontier through a transparent weighted average ($\lambda=0.25$ fixed). The adjustment is fully auditable: raw LLM probability, empirical frontier value, blending weight, and final output are directly inspectable.

\section{Methods}
\label{sec:methods}

\subsection{Empirical Frontier}

A grouped conditional logit was estimated from a Wuhan vaccination-timing DCE ($N=1{,}027$ adults, fielded December 2022--January 2023). The frontier defines $P_{\text{emp}}(x)$, the DCE-implied binary policy-versus-opt-out probability for any attribute profile $x$.

\subsection{EFR Procedure}

For each policy scenario, EFR computes:
\[
\pi_{\text{EFR}} = \lambda \cdot \pi_{\text{LLM}} + (1-\lambda) \cdot P_{\text{emp}}
\]
with $\lambda=0.25$ as an illustrative conservative operating point. The empirical component $P_{\text{emp}}$ supplies the behavioral signal; the LLM component is permitted only inside this convex blend.

\subsection{Evaluation}

Three open-source LLMs were evaluated: Qwen2.5-72B-Instruct, DeepSeek V4, and MiroThinker-1.7-mini. Two validation domains were used:

\textbf{(1) Structural consistency:} A 64-state grid spanning DCE attribute dimensions. Metrics: waiting-time monotonicity violation rate (MVR-Wait), rank correlation with $P_{\text{emp}}$.

\textbf{(2) Held-out predictive:} Six administered DCE choice tasks from 205 held-out respondents (20\% split). Metric: multinomial log loss against empirical choices.

\section{Results}
\label{sec:results}

\subsection{Held-Out Predictive Validation}

Independent local replication validation (local Ollama Qwen2.5-72B) showed multinomial log loss: Pure-DCE 0.911, EFR 0.923, Raw LLM 1.020, Uniform 1.099. EFR materially improved over Raw LLM (9.5\% reduction) and approached Pure-DCE (within 1.3\%).

\subsection{Structural Consistency}

On the 64-state grid, MVR-Wait fell from 10.4--62.5\% to 7.3--30.2\% across three LLMs under EFR anchoring (Table~1). Frontier-alignment rank correlation rose from 0.061--0.531 (Raw LLM) to 0.689--0.812 (EFR).

\section{Discussion}
\label{sec:discussion}

EFR is best understood as a constrained narrative-extension mechanism for an existing DCE. The empirical reference, not the LLM, supplies the behavioral signal. Whether narrative extension improves behavioral accuracy on out-of-design scenarios remains to be established by future work with human ground truth.

\textbf{Limitations.} The Wuhan DCE is a single-jurisdiction sample. The respondent-level holdout evaluates generalization to unseen respondents, not unseen choice tasks. Out-of-design narrative scenarios lacked human ground truth.

\textbf{Implications.} With a DCE, EFR provides a path from evidence to narrative policy simulation. Without a DCE, the framework is not a substitute for stated-preference data. Best suited for early-stage screening, not high-stakes final decisions.

\section*{References}
\addcontentsline{toc}{section}{References}

\begin{thebibliography}{10}

\bibitem{1} Aher G, et al. Using large language models for labor market text analysis. arXiv. 2022.
\bibitem{2} Chan J, et al. Discrete choice experiments in health economics. Value Health. 2024;27(1):123-135.
\bibitem{3} Hauber A, et al. ISPOR guidelines for DCEs. Value Health. 2016;19(4):405-413.

\end{thebibliography}

\end{document}
